\section{Lecture 16: Analytic Stacks and Six Functors --- Proper Maps, Open Immersions, and the Grothendieck Topology}
\label{lec:16}

Having motivated analytic stacks last time (\cref{lec:15}) as sheaves on the
opposite category of analytic rings, Scholze now turns to the objective and the
mechanism. The plan of \cref{lec:15} was: take the affine analytic spaces attached
to analytic rings as basic building blocks, and specify a Grothendieck topology on
$(\AnRing)^{\mathrm{op}}$ that tells us which gluings are allowed. The whole content
of the definition sits in \emph{which} Grothendieck topology one chooses, and the
thesis of this lecture is that this choice is dictated --- once one insists on a
good theory --- by the requirement that the assignment $A\mapsto D(A)$ underlie a
\emph{six-functor formalism}. Setting up such a formalism, extracting proper maps
and open immersions from it, and reading off the topology occupy the lecture.

Throughout ``ring'' means commutative ring and everything is (light) condensed.
Two technical caveats are deferred to the end: the set-theoretic issue that
$(\AnRing)^{\mathrm{op}}$ is a large category (so ``small sheaves'' must be used, as
in \cref{subsec:15:set-theory}), and the issue of sheaves versus hypersheaves ---
one lands on something strictly in between. We suppress both for now.

\subsection{Recollection: the same game, played before}
\label{subsec:16:recollection}

\begin{remark}[Analytic stacks by analogy with condensed sets]
\label{rmk:16:analogy}
The construction to come is very close to what was done at the very beginning of
the course. There, condensed (or light condensed) sets were defined as sheaves on
profinite sets,
\begin{equation}
\label{eq:16:condensed-as-sheaves}
\CondSetl\ =\ \mathrm{Shv}\bigl(\ProFin_{\mathbb{N}}(\mathrm{Fin})\bigr)
\ =\ \mathrm{Shv}\bigl((\text{countable Boolean algebras})^{\mathrm{op}}\bigr),
\end{equation}
starting from basic objects (profinite sets) and allowing oneself to glue them
along a specified, rather general, Grothendieck topology --- finite disjoint unions
together with \emph{all} surjections of profinite sets (\cref{def:2:light-condensed-set}).
We now play the same game with a much broader class of basic geometric objects, the
affine analytic spaces attached to analytic rings. These model many different
geometries at once --- algebraic, $p$-adic, complex analytic --- and the point of
analytic stacks is to put all of them into a single world where spaces can be built
out of pieces of different provenance. (One could equally have started from some
nicer class of basic objects, e.g.\ smooth manifolds; that too would sit inside the
world of analytic stacks.)
\end{remark}

\subsection{Desiderata: the primary invariant and the six functors}
\label{subsec:16:desiderata}

The primary invariant attached to an analytic ring is its derived category of
complete modules.

\begin{remark}[The primary invariant is a sheaf of $\infty$-categories]
\label{rmk:16:primary-invariant}
For an analytic (animated) ring $A=(\ur A,\Mod_A)$ --- most examples of interest are
static, but the general formalism is much better behaved if animated rings are
allowed --- the invariant we care about is
\begin{equation}
\label{eq:16:DA}
D(A)\ \subseteq\ D(\ur A),\qquad
D(A)\ =\ \{\,M\in D(\Cond(\ur A))\ :\ H_i(M)\in\Mod_A\ \forall i\,\},
\end{equation}
the full subcategory of the derived category of condensed $\ur A$-modules whose
homology groups are all complete (\cref{def:8:derived-cat,def:13:analytic-ring}). We
want the assignment
\begin{equation}
\label{eq:16:DA-sheaf}
A\ \longmapsto\ D(A)
\end{equation}
to be a \emph{sheaf}, and the only sensible way to phrase this is as a sheaf of
$\infty$-categories: from now on $D(A)$ is always treated as an $\infty$-category.
In particular, for any analytic stack $X$ one then \emph{defines} its
quasi-coherent derived category by descent,
\begin{equation}
\label{eq:16:DX-descent}
D(X)\ =\ \varprojlim_{\AnSpec(A)\to X}\ D(A),
\end{equation}
the limit over all affine analytic spaces mapping to $X$.
\end{remark}

We want something more than a bare sheaf of $\infty$-categories, though: we want the
whole assignment $X\mapsto D(X)$ to carry a six-functor formalism.

\begin{remark}[The six operations]
\label{rmk:16:six-functors}
A six-functor formalism organizes six operations into three adjoint pairs:
\begin{equation}
\label{eq:16:six-functors}
\otimes\ \dashv\ \RHom,\qquad f^*\ \dashv\ f_*,\qquad f_!\ \dashv\ f^!,
\end{equation}
of which the first two pairs are structure we already possess, and the third pair
is the genuinely new ingredient. Concretely:
\begin{enumerate}
\item Each $D(X)$ carries a symmetric monoidal tensor product $\otimes$, and it is
\emph{closed} --- there is an internal Hom object $\RHom$ with the usual adjunction
$\Homop(A\otimes B,C)=\Homop(A,\RHom(B,C))$.
\item For every map $f\colon Y\to X$ there is a pullback $f^*\colon D(X)\to D(Y)$,
symmetric monoidal, with right adjoint the pushforward $f_*\colon D(Y)\to D(X)$.
These come for free from the functoriality of module categories: a map of rings lets
one base-change modules, and the formalism is automatically compatible with base
change.
\item For certain maps $f\colon Y\to X$ --- the \emph{shriekable} ones --- there is a
third functor $f_!\colon D(Y)\to D(X)$, \emph{cohomology with compact support}, with
right adjoint $f^!$.
\end{enumerate}
\end{remark}

\begin{remark}[The two properties of $f_!$]
\label{rmk:16:fshriek-properties}
The functor $f_!$ is required to satisfy two properties, which are exactly what pin
down its meaning. First, \emph{base change}: for any Cartesian square
\begin{equation}
\label{eq:16:base-change-square}
\begin{tikzcd}
Y' \arrow[r, "g'"] \arrow[d, "f'"'] & Y \arrow[d, "f"] \\
X' \arrow[r, "g"'] & X
\end{tikzcd}
\end{equation}
the natural comparison $g^*f_!\xrightarrow{\ \sim\ }f'_!\,g'^*$ is an isomorphism;
in particular $f'$ is again in the class where $f_!$ is defined. Second, the
\emph{projection formula}: for $A\in D(X)$ and $B\in D(Y)$,
\begin{equation}
\label{eq:16:projection-formula}
f_!\bigl(f^*A\otimes B\bigr)\ \cong\ A\otimes f_!\,B.
\end{equation}
\end{remark}

\begin{remark}[Where this structure comes from, and coherent compact support]
\label{rmk:16:history}
This is a structure arising across mathematics. The most classical instance takes a
nice category of topological spaces (locally compact Hausdorff, say) and a suitable
derived category of sheaves; then $f_!$ exists for \emph{all} morphisms and is
literally the derived functor of sections with compact support. Historically the
formalism was first developed not there but for torsion \'etale sheaves on schemes
(Grothendieck), and later recognized for $D$-modules and other settings. One place
where it was \emph{not} much developed is the coherent (quasi-coherent) setting,
where one usually has no notion of coherent cohomology with compact support: there
is Deligne's appendix to Hartshorne's \emph{Residues and Duality}
\cite{Hartshorne_1966} sketching such a formalism, and a few scattered papers, but
it is not a theory in wide use. Our setting will naturally include a well-behaved
theory of coherent cohomology with compact support, and we very much want it.
\end{remark}

\begin{remark}[The isomorphism must be supplied: abstract six-functor formalisms]
\label{rmk:16:abstract-6ff}
As Gabber would rightly object, the description above is imprecise: when one writes
the projection-formula isomorphism \eqref{eq:16:projection-formula}, there is no
canonical comparison between the two functors --- these are \emph{a priori} two
different functors that happen to be identified, and the identification is data one
must supply. Once one starts supplying such isomorphisms, coherence problems
proliferate (base-change twice, compare the two ways, ask which compatibilities are
forced, and so on), and it was for a while unclear what the minimal way to encode all
the data in a six-functor formalism is. There is now a good $\infty$-categorical
notion of an \emph{abstract six-functor formalism}: work of Liu--Zheng
\cite{liu2024enhancedoperationsbasechange} in the setting of higher Artin stacks;
work of Gaitsgory--Rozenblyum \cite{gaitsgory2017study} in coherent cohomology (with
different names for the functors, and no talk of compact-support cohomology, but
setting up such a formalism), which is however phrased in $(\infty,2)$-categorical
language and correspondingly hard; and Mann's isolation of the key structures
\cite{mann2022padic6functorformalismrigidanalytic}, further streamlined in Scholze's
own course and notes on six-functor formalisms \cite{scholze2023six}. A moral Scholze
takes from that course:
\begin{center}
\emph{the definition of a six-functor formalism determines everything, including the
``correct'' choice of Grothendieck topology.}
\end{center}
\end{remark}

\subsection{Constructing \texorpdfstring{$f_!$}{f\textunderscore!}: proper maps, open immersions, shriekable maps}
\label{subsec:16:constructing-fshriek}

The abstract theory reduces the construction of $f_!$ to specifying two classes of
morphisms and checking a short list of conditions --- \emph{conditions}, not data.

\begin{remark}[The idea of the construction]
\label{rmk:16:fshriek-idea}
One specifies:
\begin{itemize}
\item a class of \emph{proper} morphisms, on which one \emph{declares} $f_!=f_*$; here
$f_!$ carries no new data, but one must check that this choice satisfies proper base
change and the projection formula. Since $f_!=f_*$, there is always a natural base-change
transformation and a natural projection-formula map already present, so declaring $f$
proper is just asking that these natural maps be isomorphisms --- a condition, not a datum.
\item a class of \emph{open immersions}, on which one declares $f_!$ to be the
\emph{left} adjoint $j_!$ of the pullback $j^*$; again there is a natural comparison
map in each of base change and the projection formula, and one asks that they be
isomorphisms.
\end{itemize}
The general maps for which $f_!$ is then defined --- the \emph{shriekable} maps --- are
the composites of an open immersion followed by a proper map:
\begin{equation}
\label{eq:16:shriekable-factorization}
\begin{tikzcd}[column sep=large]
Y \arrow[r, "j"] \arrow[rr, bend right=18, "f"'] & \overline Y \arrow[r, "\overline f"] & X,
\end{tikzcd}
\qquad f_!\ :=\ \overline f_*\, j_!,
\end{equation}
with $j$ an open immersion (a ``compactification'' of $Y$ inside $\overline Y$) and
$\overline f$ proper.
\end{remark}

\begin{remark}[Independence of the compactification]
\label{rmk:16:compactification-caveat}
Written as \eqref{eq:16:shriekable-factorization} this is not yet a definition: a
shriekable map admits many compactifications, and one must show $f_!$ is independent
of the choice, canonically --- and, since everything is $\infty$-categorical, not
merely up to isomorphism but with the isomorphism itself unique up to coherent higher
homotopy, and so on. This sounds like a real pain, but it is exactly what the abstract
machinery handles.
\end{remark}

\begin{theorem}[Liu--Zheng, Mann]
\label{thm:16:liu-zheng-mann}
Under very minimalistic assumptions on the classes of proper maps, open immersions,
and the resulting shriekable maps (essentially only that composites of shriekable maps
are shriekable, i.e.\ can be compactified), the above data produces an abstract
six-functor formalism. The assumptions include \emph{no} condition whatsoever on
uniqueness or domination of compactifications: for each shriekable map one need only
produce \emph{one} compactification, not compare two.
\end{theorem}

\begin{remark}[On the no-uniqueness point]
\label{rmk:16:no-uniqueness}
Gabber asked whether one must assume that two compactifications can be dominated by a
third. One does not. (Something in this direction follows for free from the mere
existence of compactifications: given two, the map to their fibre product is still a
compactification.) The one genuinely needed input is that every shriekable map is such
a composite --- open immersion followed by proper --- but only existence of a single
factorization is required. A precise statement is in Scholze's six-functor notes
\cite{scholze2023six} and in Mann's thesis \cite{mann2022padic6functorformalismrigidanalytic}.
As a sample of what the formalism is for: it should let one prove Riemann--Roch-type
statements and construct characteristic (Chern) classes in the analytic-stacks setting.
\end{remark}

\subsection{The category, and proper maps}
\label{subsec:16:proper}

We now apply these ideas to our setting.

\begin{notation}[Affine analytic spaces]
\label{not:16:aff-anstk}
Take the category
\begin{equation}
\label{eq:16:C}
\mathcal C\ =\ (\AnRing)^{\mathrm{op}}\ =\ \Aff\,\AnStk,
\end{equation}
the affine analytic (stacks) spaces. An object is written $\AnSpec(A)$; the symbol
$\AnSpec$ carries no topological space --- there is, unfortunately, no underlying
space now --- and is just a formal name for the geometric object attached to $A$,
turning $\AnRing$ into a category of geometric objects by passing to the opposite.
\end{notation}

We must decide which morphisms are proper and which are open immersions, and here one
should discard any preconceived notion of properness from algebraic geometry and read
off what the formulas demand.

\begin{definition}[Proper map]
\label{def:16:proper}
A map $f\colon\AnSpec(B)\to\AnSpec(A)$ --- i.e.\ a map of analytic rings $A\to B$ --- is
\emph{proper} if the pushforward
\begin{equation}
\label{eq:16:proper-pushforward}
f_*\colon D(B)\ \longrightarrow\ D(A),
\end{equation}
which is simply the forgetful functor (restriction of scalars along $A\to B$: any
$B$-module is an $A$-module), satisfies the projection formula.
\end{definition}

\begin{remark}[This is not the algebraic notion]
\label{rmk:16:proper-strange}
This properness is very different from the algebraic-geometry one: usually properness
is finiteness-flavoured, whereas here it will be satisfied very generally. One could
\emph{a priori} demand more of $f_*$ (proper base change, etc.), but it turns out that
once the projection formula holds, all the other good properties follow. In particular
if one embeds ordinary commutative rings into the condensed setting via the trivial
(maximal) analytic ring structure (\cref{rmk:1:trivial-structure}), then \emph{every}
map of ordinary rings is declared proper --- which looks strange, but is exactly
Huber's usage for adic spaces (see \cref{rmk:16:proper-huber}).
\end{remark}

To analyze properness, recall the two canonical operations on analytic ring
structures over a fixed base.

\begin{remark}[Induced structure and further completion]
\label{rmk:16:induced-structure}
Given a map of analytic rings $A=(\ur A,D(A))\to B=(\ur B,D(B))$, any map factors
canonically through the \emph{induced analytic ring structure} on $\ur B$: this is
the pair
\begin{equation}
\label{eq:16:induced-structure}
\bigl(\ur B,\ \{\,M\in D(\ur B)\ :\ M|_{\ur A}\in D(A)\,\}\bigr),
\end{equation}
whose complete modules are those $B^\triangleright$-modules that are complete when
restricted to $A$-modules. So an arbitrary map $A\to B$ is the composite of the
passage $A\to(\text{induced structure on }\ur B)$ followed by a \emph{further
completion} (a localization keeping the same underlying ring $\ur B$ but shrinking the
complete modules). It is precisely this further completion that corresponds
geometrically to compactification --- and the induced structures are exactly the proper
maps.
\end{remark}

\begin{proposition}[Properness $=$ inducedness]
\label{prop:16:proper-induced}
The map $f\colon\AnSpec(B)\to\AnSpec(A)$ is proper if and only if $B$ carries the
induced analytic ring structure from $A$, i.e.\ \eqref{eq:16:induced-structure} holds.
\end{proposition}

\begin{proof}
The projection formula asks that, for all $M\in D(A)$ and $N\in D(B)$,
\begin{equation}
\label{eq:16:proper-proj}
f_!\bigl(f^*M\otimes_B N\bigr)\ \cong\ M\otimes_A f_!\,N,\qquad
\text{i.e.}\quad (M\otimes_A B)\otimes_B N\ \cong\ M\otimes_A N,
\end{equation}
where on the right the tensor is taken in $D(A)$ (recall $f_!=f_*$ is the forgetful
functor). If $B$ has the induced structure, all tensor products can be computed in
$D(A)$ and the identity is clear. Conversely, take $N=\ur B$. Then the left-hand side
is the base change $M\otimes_A B$ \emph{in the sense of analytic rings}, while the
right-hand side is $M\otimes_A\ur B$ computed in $D(A)$; so the projection formula says
\begin{equation}
\label{eq:16:proper-key}
M\otimes_A B\ \cong\ M\otimes_A\ur B\qquad(\text{tensor in }D(A)).
\end{equation}
The special case $\ur A=\ur B$ is instructive: base change to $B$ is then just the
further completion, and taking $M$ complete and $N=\ur B=\ur A$ neither side does
anything ($M$ is tensored with the unit), so \eqref{eq:16:proper-key} reads $M\otimes_A
B=M$, i.e.\ the further completion is trivial and $A=B$. In general
\eqref{eq:16:proper-key} says that the base change in the sense of analytic rings
agrees with the (induced) tensor in $D(A)$ --- with no further completion needed --- which
is exactly the assertion that $B$ has the induced analytic ring structure.
\end{proof}

\begin{corollary}[Proper base change]
\label{cor:16:proper-base-change}
The class of proper maps is stable under base change, and proper base change holds.
\end{corollary}

\begin{proof}
Stability under base change is immediate from the description
\eqref{eq:16:induced-structure} (an induced structure stays induced after base change,
a straightforward unravelling of the symbols); given stability, the projection formula
after any base change gives proper base change.
\end{proof}

\begin{remark}[Comparison with Huber]
\label{rmk:16:proper-huber}
This matches Huber's notion of properness when applied to objects coming from Huber
pairs $(\ur A,A^+)$ sitting inside the analytic rings (\cref{lec:8,lec:9}). For a map
of affinoids in which the upstairs ring of integral elements $A^+$ is the smallest
possible one --- the one determined (integrally) by the $A^+$ downstairs --- Huber already
calls the map proper, and it satisfies all the good properties a proper map should. So
the seemingly strange \cref{def:16:proper} is in fact familiar from the theory of adic
spaces.
\end{remark}

\subsection{Open immersions and idempotent algebras}
\label{subsec:16:open}

For the other class one again forgets the algebro-geometric picture.

\begin{definition}[Open immersion]
\label{def:16:open}
A map $j\colon\AnSpec(B)\to\AnSpec(A)$ is an \emph{open immersion} if the pullback
$j^*$ admits a \emph{left} adjoint $j_!$ which is fully faithful (equivalently,
$j^*$ is a localization) and satisfies the projection formula. (Scholze added the
fully-faithfulness clause as an afterthought: without it one has merely a left
adjoint, not an open immersion.)
\end{definition}

\begin{remark}[The non-example]
\label{rmk:16:open-nonexample}
Any open immersion in ordinary algebraic geometry that is not also a closed immersion
is \emph{not} an open immersion in this sense. For instance
$\Gm=\AnSpec\mathbb{Z}[T^{\pm1}]\hookrightarrow\AnSpec\mathbb{Z}[T]=\mathbb{A}^1$ with
the trivial (uncompleted) ring structure: a left adjoint of $j^*$ would force pullback
$=\,\otimes\,$-of-modules to commute with all (infinite) products, and in ordinary
algebraic geometry this essentially never happens.
\end{remark}

\begin{example}[The solid $\mathbb{G}_m$]
\label{ex:16:open-example}
There is however a different route from algebraic to analytic geometry, through adic
spaces and \emph{relative solid} ring structures (\cref{lec:9}), and along it open
immersions do go to open immersions. The key example is
\begin{equation}
\label{eq:16:gm-open}
j\colon\ \AnSpec\bigl(\mathbb{Z}[T^{\pm1}]_\solid\bigr)\ \longrightarrow\ \AnSpec\bigl(\mathbb{Z}[T]_\solid\bigr),
\end{equation}
the solid $\Gm$ inside the solid affine line, which \emph{is} an open immersion. A
more primitive variant maps the solid $\Gm$ into the \emph{induced} structure on the
affine line,
\begin{equation}
\label{eq:16:gm-open-induced}
j\colon\ \AnSpec\bigl(\mathbb{Z}[T^{\pm1}]_\solid\bigr)\ \longrightarrow\ \AnSpec\bigl((\mathbb{Z}[T],\mathbb{Z})_\solid\bigr),
\end{equation}
of which \eqref{eq:16:gm-open} is essentially a base change.
\end{example}

\begin{remark}[Compact support: functions vanishing near infinity]
\label{rmk:16:compact-support-picture}
Why should such a $j_!$ deserve the name ``cohomology with compact support''?
Consider the global sections (cohomology) of the structure sheaf of the affine line,
$j_!\mathcal{O}_{\AnSpec(\mathbb{Z}[T])_\solid}$: this should be compactly supported
coherent cohomology of $\mathbb{A}^1$. A compactly supported function is one vanishing
near $\infty$. A function is an element of $\mathbb{Z}[T]$; functions near $\infty$ are
power series in $T^{-1}$, so ``vanishing at $\infty$'' means lying in the fibre of the
restriction map
\begin{equation}
\label{eq:16:vanish-at-infinity}
\mathbb{Z}[T]\ \longrightarrow\ (\text{functions near }\infty).
\end{equation}
A nonzero polynomial is determined by its germ at $\infty$, so this map is
\emph{injective} --- the naive kernel vanishes --- but as a \emph{complex} the fibre is
still interesting: its cokernel is nontrivial, in fact a product of copies of
$\mathbb{Z}$ (the coefficients of the strictly-negative powers of $T$). Thus
$j_!\mathcal{O}$ is a genuine two-term complex, exhibiting compact-support cohomology
even though no honest compactly-supported functions exist.
\end{remark}

Concretely, the completion $j_*j^*$ is computed by an internal Hom out of the object
$C'=j_!\mathbb{Z}[T]$, recovering the solidification formula of \cref{lec:7,lec:9}.

\begin{proposition}[The completion $j_*j^*$]
\label{prop:16:jshriek-formula}
For the open immersion of \cref{ex:16:open-example}, and $M\in D(\mathbb{Z}[T]_\solid)$,
\begin{equation}
\label{eq:16:jshriek-rhom}
j_*\,j^*M\ =\ \RHom_{\mathbb{Z}[T]}\bigl(C',\,M\bigr),\qquad
C'\ :=\ j_!\,\mathbb{Z}[T]\ =\ \operatorname{fib}\bigl(\mathbb{Z}[T]\hookrightarrow\mathbb{Z}(\!(T^{-1})\!)\bigr),
\end{equation}
the two-term complex $\bigl[\,\mathbb{Z}[T]\to\mathbb{Z}(\!(T^{-1})\!)\,\bigr]$ (with
$\mathbb{Z}[T]$ in homological degree $0$ and the differential the inclusion into the
``functions near $\infty$'' algebra), whose only nonzero homology is
$H_{-1}(C')=\mathbb{Z}(\!(T^{-1})\!)/\mathbb{Z}[T]\cong\prod_{\mathbb{N}}\mathbb{Z}$; in
particular
\begin{equation}
\label{eq:16:jshriek-O}
j_!\,\mathcal{O}_{\AnSpec(\mathbb{Z}[T])_\solid}\ =\ j_!\,\mathbb{Z}[T]\ =\ C'.
\end{equation}
This is precisely the $\mathbb{Z}[T]$-solidification formula Clausen gave for the
complement of the closed disk (\cref{prop:7:solidification-formula},
\cref{eq:9:complement-rhom}).
\end{proposition}

Reading \eqref{eq:16:jshriek-rhom} the other way, tensoring against the relevant object
gives the left adjoint $j_!$ as a $\otimes$-with-a-module, which is exactly what one
needs to verify the projection formula.

\begin{proposition}[Open immersions $=$ idempotent algebras]
\label{prop:16:open-idempotent}
Let $X=\AnSpec(A)$. The open immersions $j\colon Y=\AnSpec(B)\to X$ are equivalent to
idempotent commutative algebra objects $C\in D(A)$ such that $\RHom_A(C,-)[1]$ commutes
with all limits and preserves connectivity ($D_{\ge0}(A)$), via
\begin{equation}
\label{eq:16:open-to-idempotent}
j\ \longmapsto\ C\ =\ \operatorname{cone}\bigl(j_!\mathbf 1\to\mathbf 1\bigr).
\end{equation}
\end{proposition}

\begin{proof}[Sketch]
Since a right adjoint to $j_!$ automatically exists (namely $j^*$, once $j_!$ is fully
faithful --- there is a localization here), one recovers $B$ as the reflective
subcategory
\begin{equation}
\label{eq:16:killed-by-C}
D(B)\ =\ \{\,M\in D(A)\ :\ \RHom_A(C,M)=0\,\},
\end{equation}
the objects killed by $C$ in the sense of \cref{def:13:killing}. To produce an
idempotent algebra one needs only supply the unit $\mathbf 1\to C$ and check that
$C\otimes_A(\mathbf 1\to C)$ becomes an isomorphism; the projection formula for $j_!$
does the rest. Indeed $j_!\mathbf 1$ is itself an idempotent (co)algebra:
\begin{equation}
\label{eq:16:jshriek-idempotent}
j_!\mathbf 1\otimes_A j_!\mathbf 1\ =\ j_!\bigl(j^* j_!\mathbf 1\bigr)\ =\ j_!\mathbf 1,
\end{equation}
so its cone $C$ is an idempotent algebra, describing a ``space near infinity'' --- the
complementary closed subset cut out by the near-infinity functions. The two conditions
on $C$ (that the completion $\RHom_A(C,-)$ commute with all limits and preserve
connectivity) are exactly what is needed for this reflective localization to define an
analytic ring structure; all other analytic-ring axioms follow formally
(\cref{def:13:killing,prop:13:left-adjoint}).
\end{proof}

\begin{remark}[$j_!j^*$ via the projection formula]
\label{rmk:16:jshriek-proj}
The projection formula gives $j_!j^*M=j_!\mathbf 1\otimes_A M$, and dually
$j_*j^*M=\RHom_A(j_!\mathbf 1,M)$. So the whole new (localized) analytic ring structure
is determined by the single object $j_!\mathbf 1$, or equivalently by $C$ (recover
$j_!\mathbf 1$ as $\operatorname{fib}(\mathbf 1\to C)$).
\end{remark}

\begin{corollary}[Base change]
\label{cor:16:open-base-change}
The class of open immersions is stable under base change, and the functors $j_!$, $j^*$
commute with base change.
\end{corollary}

\subsection{Shriekable maps and the six-functor formalism}
\label{subsec:16:shriekable}

With proper maps and open immersions in hand, the shriekable maps are their
composites --- and here the factorization is canonical, because
\cref{prop:16:proper-induced} already tells us which proper map to take.

\begin{definition}[Shriekable map]
\label{def:16:shriekable}
A map $f\colon\AnSpec(B)\to\AnSpec(A)$ is \emph{shriekable} if, in its canonical
factorization through the induced structure,
\begin{equation}
\label{eq:16:canonical-factorization}
\begin{tikzcd}[column sep=large, row sep=large]
\AnSpec(B) \arrow[r, "j"] \arrow[rr, bend right=16, "f"'] & \AnSpec\bigl(A\otimes_{\ur A}\ur B\bigr) \arrow[r, "\overline f"] & \AnSpec(A),
\end{tikzcd}
\end{equation}
the first map $j$ --- from $B$ to the induced analytic ring structure on $\ur B$ --- is an
open immersion. (The second map $\overline f$ is proper by
\cref{prop:16:proper-induced}, and $j$ automatically has $j_!$ shriekable-friendly; the
only real condition is that $j$ be an open immersion.) One then sets
$f_!=\overline f_*\,j_!$.
\end{definition}

\begin{example}[Structure map of the solid $\Gm$]
\label{ex:16:shriekable-gm}
Taking the solid structure on $\mathbb{Z}[T^{\pm1}]$, the structure map
\begin{equation}
\label{eq:16:gm-shriekable}
f\colon\ \AnSpec\bigl(\mathbb{Z}[T^{\pm1}]_\solid\bigr)\ \longrightarrow\ \AnSpec(\mathbb{Z}_\solid)
\end{equation}
is shriekable: the compactification $\overline f$ is exactly the induced structure, and
$j$ is the open immersion of \cref{ex:16:open-example}. More generally, any map of
finite-type algebras equipped with the relative solid ring structure gives shriekable
maps.
\end{example}

\begin{remark}[Huber's ``$+$-weakly finite type'']
\label{rmk:16:weakly-ft}
For a map of Huber pairs, being shriekable is exactly the condition Huber calls being
of \emph{$+$-weakly finite type}: that the upstairs ring of integral elements $A^+$ is
generated, as a ring of integral elements, by finitely many new elements over the
base. Huber introduced this class of morphisms (and the associated proper/open
notions) precisely as the maps for which one can define compact-support functors in
\emph{torsion \'etale} cohomology of adic spaces; remarkably, they are now also the
maps for which one gets compact-support functors in \emph{coherent} cohomology.
\end{remark}

\begin{proposition}[The six-functor formalism on affine analytic spaces]
\label{prop:16:6ff-exists}
The data of proper maps (\cref{def:16:proper}), open immersions (\cref{def:16:open}) and
the resulting shriekable maps satisfies all the conditions required by
\cref{thm:16:liu-zheng-mann}, and hence produces a six-functor formalism
\begin{equation}
\label{eq:16:6ff-on-C}
X\ \longmapsto\ D(X)\qquad\text{on}\qquad \mathcal C\ =\ (\AnRing)^{\mathrm{op}}\ =\ \Aff\,\AnStk,\ \text{with shriekable maps.}
\end{equation}
\end{proposition}

\begin{remark}[The small checks]
\label{rmk:16:small-checks}
Several small verifications are suppressed (each an easy check): that a composite of two
shriekable maps is again shriekable --- essentially by base change --- and one
compatibility between the lower-shriek and the lower-star ($=$ lower-shriek) functors of
the proper maps. Granting these, \cref{prop:16:6ff-exists} is a corollary of the general
existence result. This is what was promised, but not fully constructible, in the first
Clausen--Scholze course on condensed mathematics: there one specialized to adic spaces
with underlying ring a field and solid modules, said such a formalism should exist, and
now it does --- and from it proofs of Poincar\'e duality, Serre duality, and the like
follow easily.
\end{remark}

\subsection{From affine analytic spaces to analytic stacks}
\label{subsec:16:topology}

So far we have a six-functor formalism on the affine objects. To reach genuine
geometric objects --- built by gluing affines, as schemes are built from affines and
stacks from schemes --- one wants to extend it to sheaves, and the extension itself
selects the Grothendieck topology.

\begin{remark}[The general extension question]
\label{rmk:16:extension-question}
Given a six-functor formalism $X\mapsto D(X)$, a functor $\mathcal C\to\Cat$ on some
$\infty$-category $\mathcal C$, one wants to extend it (for some notion of sheaves) to
the sheaves on $\mathcal C$,
\begin{equation}
\label{eq:16:extend-to-sheaves}
\mathrm{Shv}\bigl(\mathcal C,\ \mathrm{Anima}\bigr),
\end{equation}
i.e.\ to stacks glued out of objects of $\mathcal C$, on which the extended $D$ becomes a
sheaf of $\infty$-categories. This was in fact the original motivation of Liu--Zheng
\cite{liu2024enhancedoperationsbasechange}, who wanted to extend torsion \'etale
cohomology from schemes to (higher) Artin stacks. Scholze's six-functor notes
\cite{scholze2023six} analyze this question --- following Mann's rephrasing
\cite{mann2022padic6functorformalismrigidanalytic}, then rephrased again --- and try to
pin down the best Grothendieck topology, including a general definition of what Scholze
there calls the ``$\mathsf D$-topology''. He is not entirely happy with the precise
combinatorics in full generality, but the takeaway is clean.
\end{remark}

\begin{remark}[The takeaway: universal descent]
\label{rmk:16:takeaway}
The covers should be exactly those that satisfy \emph{universal $*$-descent} and
\emph{universal $!$-descent}. Precisely:
\begin{enumerate}
\item A map $f\colon Y\to X$ satisfies \emph{$*$-descent} if the pullback exhibits
$D(X)$ as the totalization of the \v Cech diagram of $f$,
\begin{equation}
\label{eq:16:star-descent}
D(X)\ \xrightarrow{\ p^*\ }\ \varprojlim_{\Delta}\ \Bigl(\ D(Y)\ \rightrightarrows\ D(Y\times_X Y)\ \Rrightarrow\ \cdots\ \Bigr),
\end{equation}
all transition functors being upper-stars $T_1^*,T_2^*,\dots$; it satisfies
\emph{universal $*$-descent} if this holds after any base change.
\item A shriekable map $f\colon Y\to X$ satisfies \emph{$!$-descent} if the same holds
with the upper-shriek functors $f^!$ in place of the upper-stars, and \emph{universal
$!$-descent} if this persists after any base change.
\end{enumerate}
Here ``any base change'' should be read as base change staying within $\mathcal C$; the
fully general (base change to any presheaf) formulation is more awkward and is what
makes the abstract topology delicate.
\end{remark}

\begin{theorem}[$!$-descent suffices]
\label{thm:16:descent}
If a shriekable map $f\colon Y\to X$ in $\Aff\,\AnStk$ satisfies $!$-descent, then it
satisfies both universal $!$-descent and universal $*$-descent. In fact it satisfies
something stronger: descent even at the level of the $(\infty,2)$-category of module
categories --- the assignment $A\mapsto\mathrm{Pr}^L_{D(A)}$ of presentable stable
$\infty$-categories linear over $D(A)$ satisfies descent. Asking for $!$-descent is
equivalent to asking descent at this two-categorical level (which is in fact how the
theorem is proved).
\end{theorem}

\begin{definition}[The analytic-stacks topology]
\label{def:16:topology}
Endow $\Aff\,\AnStk$ with the Grothendieck topology generated by:
\begin{enumerate}
\item finite disjoint unions --- whenever $I$ is finite and the $X_i$ are affine analytic
spaces, $\coprod_{i\in I}X_i$ is covered by the individual $X_i$ (as for profinite sets,
where one allowed finite disjoint unions); and
\item shriekable maps satisfying $!$-descent.
\end{enumerate}
Analytic stacks are then (a suitable class of hyper-)sheaves on $\Aff\,\AnStk$ for this
topology.
\end{definition}

\begin{remark}[Descendability, $h$-covers, and how general this is]
\label{rmk:16:descendability}
This is a very general class. Restricting to proper maps, $!$-descent (there $!=*$) is
precisely the condition that the map of algebras be \emph{descendable} in the sense of
Mathew \cite{mathew2016galoisgroupstablehomotopy} --- a robust, higher-categorical notion
of descent satisfied by virtually all faithfully flat maps that are countably presented
(one reason to work in the light condensed setting: light rings are countably
presented). Descendable maps form a broad class: they include all countably presented
faithfully flat maps, all quotients by nilpotent ideals, and --- more surprisingly --- all
$h$-covers in Voevodsky's $h$-topology. Over Noetherian affine schemes, for a
finitely-presented map of Noetherian affines, descendability is equivalent to being an
$h$-cover. (For usual commutative rings with the trivial structure one declares all maps
of affines proper, which looks awkward but works: then all maps are shriekable, and
$!$-descent picks out exactly the descendable, i.e.\ $h$-covering, maps.) For the
$(\infty,2)$-categorical descent statements one should consult Mathew's paper.
\end{remark}

\begin{remark}[The two deferred caveats]
\label{rmk:16:caveats}
As flagged at the start, two technical points remain. First, set theory: as in
\cref{subsec:15:set-theory} one must work with accessible presheaves, and check that the
above topology has good approximation properties, i.e.\ that sheafification preserves
accessibility. Second, hyper-completeness: as already in the light condensed setting one
does not want plain sheaves but a certain class of hypersheaves, so
\cref{def:16:topology} really allows a controlled class of hypercovers --- the same
maneuver as before. Both are suppressed here.
\end{remark}

\begin{remark}[What ``module categories'' descent means]
\label{rmk:16:2cat-descent}
The strong form of \cref{thm:16:descent} uses Lurie's category $\mathrm{Pr}^L$ of
presentable $\infty$-categories and colimit-preserving functors \cite{lurie2017higher};
inside it, the stable ones linear over a base $D(A)$ form the objects
$\mathrm{Pr}^L_{D(A)}$, morally $D(A)$-linear DG-categories (e.g.\ $D(B)$ for a
$D(A)$-algebra $B$). One asks whether $A\mapsto\mathrm{Pr}^L_{D(A)}$ satisfies descent
(a $*$-type descent: base-change module categories, ask the presheaf to be a sheaf).
This is naturally an $(\infty,2)$-category, but for the descent question one may forget
the non-invertible $2$-morphisms and work with the underlying $(\infty,1)$-category ---
those extra $2$-morphisms automatically satisfy descent as well. In this sense the
$(\infty,2)$-categorical structure does not matter much, and $2$-categorical
$*$-descent is equivalent to $1$-categorical $!$-descent one level down.
\end{remark}
